\section{Limitations and the status of each number}
\label{sec:limitations}

We separate three kinds of evidence, because conflating them is the easiest way
to overstate a study of this shape.

\paragraph{Run directories are not independent observations.}
This project produced \nRunDirs{} completed full-budget run directories, holding
\nDistinctRuns{} distinct training trajectories: \nDuplicateGroups{} groups
covering \nDuplicateDirs{} historical directories are bitwise identical to one
another because Waves~7--10 re-ran the same scientific configuration to measure
systems changes (Table~\ref{tab:evidence}), whereas all four runs in
\texttt{wave\_sym\_fastlag} are distinct trajectories. We verified this by
comparing full 5{,}120-point training-loss series rather than trusting wave
labels, and every count in this paper counts each bitwise-identical group once.

\begin{table}[t]
\centering
\caption{Bitwise-identical run groups across Waves~7--10. Each group is one
trajectory recorded under several wave labels while systems changes were
measured around it; the recorded losses agree at all 5{,}120 updates.}
\label{tab:evidence}
\footnotesize\setlength{\tabcolsep}{6.0pt}
% Generated by scripts/check_duplicates.py. Do not edit.
\begin{tabular}{lrl}
\toprule
Trajectory & Directories & Appears under \\
\midrule
\texttt{wave10/muont} & 4 & wave10, wave7, wave8, wave9 \\
\texttt{wave10/outer\_leftdiag\_fused} & 3 & wave10, wave8, wave9 \\
\texttt{wave10/null\_leftdiag\_upstream\_fused} & 2 & wave10, wave9 \\
\texttt{wave10/outer\_leftdiag\_upstream\_fused} & 2 & wave10, wave9 \\
\texttt{wave7/outer\_leftdiag} & 2 & wave7, wave8 \\
\bottomrule
\end{tabular}

\end{table}

\paragraph{Which claims are replicated, and which are single-seed.}
The complementarity result (Section~\ref{sec:complement}), the freshness result
(Section~\ref{sec:freshness}), and the symmetric width-2048 ablation and
eight-A100 \texttt{fast\_lag} wall-clock crossings
(Sections~\ref{sec:sym-ablation} and~\ref{sec:systems}) rest on seed-0
comparisons with exact token pairing and standard errors two orders of magnitude
below the paired training effects. \taxid{WT-MLP}'s quality benefit is
replicated on three seeds (Section~\ref{sec:replication}), with a factor-of-two
spread and the largest value on the selection seed, while the matched-null
effect is observed in three separate nulls across Waves~1, 3, and~10. Nothing
here is replicated on a seed reserved from exploration: seeds 0, 1 and 2 have
all been inspected, which is why Section~\ref{sec:program} reserves 3, 4 and 5.

\paragraph{Known confounds and what remains open.}
Three boundaries govern how the current numbers should be read.
\emph{Full-budget learning-rate neighborhood}: while Section~\ref{sec:tuning}
shows that the 320-step optimum ($0.064$) regresses full-budget \muont{} by
\waveTwelveMuontSelDiffMag{} nats and Section~\ref{sec:mech-graft} rules out
step-size inflation (\mechMlpdiagGraftNorm{} norm ratio), the $\{0.016, 0.032,
0.045\}$ full-budget points around our $0.0226$ anchor remain to be executed
(Section~\ref{sec:program}, Item~1).
\emph{Same-wave \texttt{fast\_lag} matched null and baseline polling jitter}:
Wave-3 \nullmlp{} crossed a wave boundary relative to Wave-2 \mlpdiag{}, and
\texttt{wave\_sym\_fastlag} has not yet executed its own \texttt{fast\_lag}
\taxid{Null-WT} control (Item~2; Table~\ref{tab:main} marks net ${}^{\ddagger}$
deltas accordingly). In addition, \texttt{wave\_sym\_fastlag/muont} is bitwise
identical to \texttt{wave0/muont} through update~\symWaveExactSteps{} and
diverges by $\symWaveMuontVsWaveZeroSel{}$ late selection nats thereafter due to
upstream asynchronous host \texttt{eigh} polling
(Appendix~\ref{app:systems}); we therefore report \texttt{wave\_sym\_fastlag}
deltas against both same-wave \taxid{MuonT} and \texttt{wave0/muont}.
\emph{Single-seed wall-clock margin}: \taxid{WT-Sym-Lag}'s launch-anchored lead
at $\text{C4}\le 3.52$ ($\sysCrossSymLagThreeFiveTwoS$\,s versus
$\sysCrossMuonWaveZeroThreeFiveTwoS$\,s for \taxid{Muon} and
$\sysCrossMuontSymThreeFiveTwoS$\,s for same-wave \taxid{MuonT}) is a modest
$8$--$107$\,s margin on one seed and requires multi-seed confirmation in Item~4.

\paragraph{Reproduction fidelity.}
Our baseline is a runtime adaptation, not an exact reproduction of the pinned
reference: A100 with FlashAttention-2 rather than H100 with FA3, microbatch 4
with 4-way accumulation rather than 8 with 2-way, deterministic algorithms
enabled, and explicit compiler options rather than the reference mode argument.
The data and budget are mathematically identical. Our Wave-0 baseline's final C4
sits $+0.0021$ nats from the pinned reference value while plain \taxid{Muon} sits
$-0.0003$ from its own, so our reproduction of the baseline-to-preconditioned
gap is $0.0091$ against a reference $0.0115$ (and $0.0107$ in
\texttt{wave\_sym\_fastlag}). That residual ($0.0008$--$0.0024$ nats) does not
invalidate paired within-wave comparisons, which share every implementation
detail, but it bounds how much weight any single final evaluation can carry on
its own.

\paragraph{Generality.}
One model size, one token budget, one hardware configuration, one data mixture.
We have one larger-batch pair and no larger model. The arm geometries are tied
to this architecture's projection structure --- in particular to the SwiGLU MLP
expansion ratio of 4, which concentrates $96.2\%$ of $O(d^3)$ eigendecomposition
work in the width-2048 projections.

\section{Related work}
\label{sec:related}

\paragraph{Muon and matrix-aware steps.}
Muon \citep{jordan2024muon} established that orthogonalizing momentum with a
cheap fixed-degree iteration is competitive at scale, and the preconditioned
variant we baseline against adds input-side whitening to it. Our arms are
modifications of that baseline, not of Muon itself, and we report plain Muon
only as a reference point. We are not aware of a published study that adds
output-cotangent second moments to a Muon-type optimizer, separates pre- versus
post-SwiGLU width-2048 covariance regimes, and measures the result against a
matched execution-path null; where our arms approach a published two-sided form
we label them formula controls rather than reproductions.

\paragraph{Kronecker-factored and full-matrix preconditioners.}
K-FAC \citep{martens2015kfac} and the Shampoo family
\citep{gupta2018shampoo,anil2020scalable} use both sides of the gradient, and
our motivation is theirs. Two differences shape our results. First, those
methods do not orthogonalize, so their preconditioner supplies conditioning
that Newton--Schulz already supplies here --- consistent with our finding that
at width 2048 on pre-SwiGLU outputs, current diagonal scale beats a stale
eigenbasis, whereas on post-SwiGLU \texttt{in:down} off-diagonal structure still
helps during early constant-LR training. Second, Shampoo-family implementations
routinely accept preconditioner staleness of hundreds of steps, whereas we find
a four-update eigenbasis staleness already costly on rapidly drifting cotangents
while a one-update factor lag ($\tau=1$, one-sixteenth of the EMA window) hides
almost all side-stream latency. SOAP \citep{vyas2024soap} occupies the middle
ground and is a natural comparison under a tuned protocol.

\paragraph{Empirical Fisher and what $B$ measures.}
Our $B$ is an uncentered moment of observed cotangents, the quantity whose gap
from the true Fisher is analyzed by \citet{kunstner2019limitations}. We take
their caution seriously enough not to call our construction curvature-based:
Section~\ref{sec:program} specifies the factorization check that would be
needed to earn that description, and it has not been run.

\paragraph{Optimizer evaluation methodology.}
That optimizer comparisons are dominated by tuning protocol has been argued
forcefully \citep{schmidt2021crowded,choi2019empirical,kaddour2023nodata}. Our
contribution is a concrete instance with the horizon mechanism exposed: a
320-step screening proxy selected an over-stepped learning rate ($0.064$) that
degraded the 5{,}120-step baseline itself by \waveTwelveMuontSelDiffMag{} nats.
We add the matched execution-path null as a second discipline that accounted for
a fifth to a third of raw fixed-rate gains.

\paragraph{Newton--Schulz as a finite map.}
Our step-budget and operand-precision measurements (Appendix~\ref{app:numerics})
treat the iteration as the finite, rounded object it is, and find a direction
plateau at seven steps. Related work on compensated and mixed-precision matrix
products \citep{henry2019leveraging} informs our Ozaki FP16-split Tensor-Core
update stack in Section~\ref{sec:systems}.
